\section{Edge contractions, residues, and bar signs}
\label{sec:edge-contraction-signs}

A two-step collision has two orders. Its two contributions cancel only when
the signs oppose and the resulting coefficient maps satisfy the required identity.
The determinant of the remaining edges records the first condition.
An associative product supplies the second condition in the ordinary bar construction.
Logarithmic residues give a geometric realization of the same determinant sign.

All complexes use cohomological grading over a commutative ring $k$.
Thus $V[r]^j=V^{j+r}$ and $d_{V[r]}=(-1)^r d_V$.
Tensor products use the Koszul rule.
Every sum below is algebraic and finite on each input.
No completion or passage to graph isomorphism classes is implicit.

\subsection{The deciding corner}
\label{subsec:deciding-corner}

Let $e<f$ be two labeled edges and let $\omega=e\wedge f$.
Place this determinant generator in degree $-2$.
Deleting an edge means exterior contraction with its dual basis vector:
\begin{equation}
\label{eq:two-edge-deletion}
\iota_e\omega=f,\qquad \iota_f\omega=-e,\qquad
\iota_f\iota_e\omega=1,\qquad \iota_e\iota_f\omega=-1.
\end{equation}
Suppose contracting $e$ then $f$ gives a coefficient map $P$.
Suppose the reverse order gives $Q$, with the same source and target.
The coefficient of the double contraction is $P-Q$.
For $P=Q$ it vanishes over every $k$.
For $k=\mathbb Z$, $P=1$, and $Q=2$ on $\mathbb Z$, it equals $-1$.
The determinant sign therefore leaves a coefficient identity to prove.

The same calculation appears on a coordinate bidisc over $\mathbb C$.
Write $\lambda_e=du/u$ and $\lambda_f=dv/v$.
Normalize the residue by placing the normal logarithmic form first:
\begin{equation}
\label{eq:left-residue}
r_u(\lambda_e\wedge\beta+\gamma)=\beta|_{u=0},
\end{equation}
where $\gamma$ has no logarithmic pole along $u=0$.
Then
\[
r_u(\lambda_e\wedge\lambda_f)=\lambda_f,\qquad
r_v(\lambda_e\wedge\lambda_f)=-\lambda_e.
\]
The two ordered residues are $1$ and $-1$.
There is already one determinant sign in this formula.
Multiplying each residue by its deletion sign again gives $1+1=2$.
This last prescription fails over $\mathbb Z$ and $\mathbb C$.

\subsection{A determinant complex with two contractions}
\label{subsec:determinant-complex}

Fix a finite set $E$ of labeled edges.
For $S\subseteq E$, write $F=E\setminus S$ and set
\[
L(F)=\bigwedge^{|F|}k^F[|F|].
\]
The unshifted exterior power lies in degree zero, so $L(F)$ lies in degree $-|F|$.
Its differential is zero.
For $e\in F$, contraction defines a degree-one map
$\iota_e:L(F)\longrightarrow L(F\setminus\{e\})$.
An ordering $F=(e_1,\ldots,e_r)$ gives
\begin{equation}
\label{eq:ordered-deletion}
\iota_{e_i}(e_1\wedge\cdots\wedge e_r)
=(-1)^{i-1}e_1\wedge\cdots\widehat e_i\cdots\wedge e_r.
\end{equation}
This formula is compatible with every relabeling of $F$.
For $i<j$, the signs for deleting $e_i$ then $e_j$, and conversely, are
$(-1)^{i+j-3}$ and $(-1)^{i+j-2}$.
Hence $\iota_f\iota_e=-\iota_e\iota_f$ for distinct edges.
Once $e$ is removed, a second contraction of $e$ has no domain.
There is no repeated-edge summand.

For every $S\subseteq E$, let $(M_S,d_S)$ be a cochain complex of $k$-modules.
For $e\notin S$, supply two degree-zero cochain maps
\[
a_{S,e},b_{S,e}:M_S\longrightarrow M_{S\cup\{e\}}.
\]
The target records the contracted set, so both contraction orders have the same target.
The coefficient maps can be zero.
Define
\[
C=\bigoplus_{S\subseteq E} L(E\setminus S)\otimes M_S,
\qquad |\omega\otimes x|=|x|-|E\setminus S|.
\]
The determinant factor comes first.
For $r=|E\setminus S|$, put
\begin{align}
\label{eq:total-contractions}
\delta(\omega\otimes x)&=(-1)^r\omega\otimes d_Sx,\notag\\
A(\omega\otimes x)&=\sum_{e\notin S}\iota_e\omega\otimes a_{S,e}x,\\
B(\omega\otimes x)&=\sum_{e\notin S}\iota_e\omega\otimes b_{S,e}x.\notag
\end{align}
All three maps have degree one.

For distinct $e,f\notin S$, abbreviate $S\cup\{e\}$ by $S+e$ and define
\begin{align}
\label{eq:coefficient-defects}
K^{aa}_{S;e,f}
 &=a_{S+e,f}a_{S,e}-a_{S+f,e}a_{S,f},\notag\\
K^{bb}_{S;e,f}
 &=b_{S+e,f}b_{S,e}-b_{S+f,e}b_{S,f},\\
K^{ab}_{S;e,f}
 &=a_{S+e,f}b_{S,e}+b_{S+e,f}a_{S,e}\notag\\
 &\quad-a_{S+f,e}b_{S,f}-b_{S+f,e}a_{S,f}.\notag
\end{align}
These are degree-zero maps $M_S\to M_{S+e+f}$.

\begin{theorem}[Exact coefficient criterion]
\label{thm:exact-coefficient-criterion}
The operators satisfy $\delta^2=\delta A+A\delta=\delta B+B\delta=0$.
Moreover,
\[
A^2=0\iff K^{aa}_{S;e,f}=0,\qquad
B^2=0\iff K^{bb}_{S;e,f}=0,
\]
and
\begin{equation}
\label{eq:exact-mixed-criterion}
AB+BA=0\iff K^{ab}_{S;e,f}=0
\quad\text{for every }S,e,f.
\end{equation}
Consequently $\delta+A+B$ is a differential when all three coefficient identities hold.
These statements hold in characteristic two as well.
\end{theorem}

\begin{proof}
The assertion $\delta^2=0$ follows from $d_S^2=0$.
For a fixed edge $e$, the two terms in $\delta A+A\delta$ have signs
$(-1)^{r-1}$ and $(-1)^r$.
Their coefficient maps agree because $d_{S+e}a_{S,e}=a_{S,e}d_S$.
The proof for $B$ uses the corresponding cochain-map identity.

Fix an ordering of $E\setminus S$.
For $e=e_i<f=e_j$, the component of $A^2$ from $S$ to $S+e+f$ equals
\[
(-1)^{i+j-3}\omega_{F\setminus\{e,f\}}
 \otimes K^{aa}_{S;e,f}x.
\]
There are exactly two contributing paths.
Equation~\eqref{eq:ordered-deletion} gives their signs.
The same calculation gives $K^{bb}$ for $B^2$.
There are four paths in $AB+BA$.
Two first remove $e$, and two first remove $f$.
Their sum has the same determinant factor and coefficient $K^{ab}$.
Distinct unordered pairs have distinct target summands for a fixed $S$.
The determinant factor is free of rank one.
Thus vanishing of each operator is equivalent to the displayed coefficient identities.
No step divides by two.
\end{proof}

\begin{corollary}[Fixed edge colors]
\label{cor:colored-edges}
Partition $E=E_a\sqcup E_b$.
Suppose degree-zero cochain maps $m_{S,e}$ satisfy
\[
m_{S+e,f}m_{S,e}=m_{S+f,e}m_{S,f}.
\]
Set $a_{S,e}=m_{S,e}$ on $E_a$ and zero on $E_b$.
Set $b_{S,e}=m_{S,e}$ on $E_b$ and zero on $E_a$.
Then $A^2=B^2=AB+BA=0$.
\end{corollary}

\begin{proof}
For two edges of one color, the pure coefficient defect is the assumed square.
For different colors, the mixed defect is that square.
All other terms are zero.
\end{proof}

When both operations can contract every edge, the four-term identity is the exact condition.
Separate associativity of the two operations does not imply it.
For a tree with a fixed edge labeling, its quotients by subsets give the required target identifications.
Coefficient complexes and contraction maps must also respect these identifications.
If relabelings act compatibly, the construction descends to their coinvariants by equivariance.
The vanishing criterion on a quotient can be weaker, since distinct labeled terms can become equal there.

\subsection{Stable graphs with labeled edges}
\label{subsec:stable-graph-contractions}

A stable graph here is a connected graph with finitely many vertices, edges, and labeled legs.
Each vertex $v$ has a nonnegative integer genus $g_v$ and satisfies
$2g_v-2+\operatorname{val}(v)>0$.
A loop contributes two to the valence.
For a connected graph, its total genus is $\sum_v g_v+|E|-|V|+1$.
Fix an edge labeling and, when needed, an edge color.

For $S\subseteq E$, the vertices of $\Gamma/S$ are the connected components
$C$ of the subgraph $(V,S)$, including isolated vertices.
Give the vertex corresponding to $C$ genus
\begin{equation}
\label{eq:simultaneous-graph-genus}
g_C=\sum_{v\in C}g_v+|S_C|-|V_C|+1.
\end{equation}
Retain the edges outside $S$, with endpoints sent to their components.
Retain each leg and its label.
This defines the quotient without choosing a contraction order.
A nonloop edge merges two vertices and adds their genera.
A loop increases its vertex genus by one.

\begin{proposition}[The oriented stable-graph differential]
\label{prop:stable-graph-signs}
The quotient $\Gamma/S$ is stable and preserves total genus.
For disjoint $S,T\subseteq E$, contracting $S$ and then $T$ gives $\Gamma/(S\cup T)$.
On the free module of labeled stable graphs with determinant factor $L(E)$,
edge contraction defines a differential of cohomological degree one.
Splitting the edges into two fixed colors gives anticommuting contraction differentials.
The operators descend to graph-isomorphism coinvariants if isomorphisms preserve leg labels and edge colors.
\end{proposition}

\begin{proof}
The number $|S_C|-|V_C|+1$ is the cycle rank of the connected subgraph $C$, so $g_C\geq0$.
Its surviving valence is $\sum_{v\in C}\operatorname{val}(v)-2|S_C|$.
Therefore
\[
2g_C-2+\operatorname{val}(C)
=\sum_{v\in C}\bigl(2g_v-2+\operatorname{val}(v)\bigr)>0.
\]
Summing~\eqref{eq:simultaneous-graph-genus} over $C$ and then adding the cycle rank of $\Gamma/S$
returns $\sum_vg_v+|E|-|V|+1$.
To compare successive quotients, fix one component of $(V,S\cup T)$.
If it contains $c$ components after contracting $S$, the second contraction contributes $|T_C|-c+1$ to its genus.
The first contraction has contributed $|S_C|-|V_C|+c$.
Their sum is $|S_C|+|T_C|-|V_C|+1$.
The vertex genera, remaining flags, and leg labels therefore agree with the simultaneous quotient.

Each two-edge coefficient square now has equal composites, because both yield that same quotient graph.
Its determinant signs are opposite by~\eqref{eq:ordered-deletion}.
This proves the square-zero and mixed identities.
Relabeling edges transports their determinant generators and contraction maps together.
Thus these operators are equivariant and descend to coinvariants.
No division by an automorphism order is used.
\end{proof}

\begin{example}[Two parallel edges]
\label{ex:parallel-edges}
Take two genus-zero vertices, each carrying one leg, joined by two edges $e<f$.
Each vertex has valence three.
Contracting $e$ gives one genus-zero vertex with two legs and the remaining loop $f$.
Contracting that loop gives one genus-one vertex with two legs.
The other order gives the same result.
The two determinant signs are $+1$ and $-1$, so the sum is zero.
The first edge is nonseparating but is not a loop.
Its contraction does not increase the vertex genus.
\end{example}

\subsection{An explicit comparison with the associative bar construction}
\label{subsec:associative-comparison}

Let $(H,d_H,\mu)$ be an augmented differential graded associative $k$-algebra.
Write $\bar H$ for its augmentation ideal.
Let $t$ have degree $-1$, and put $t\bar H=\bar H[1]$.
The reduced tensor coalgebra is
\[
T^c(t\bar H)=\bigoplus_{n\geq0}(t\bar H)^{\otimes n},
\]
with deconcatenation coproduct and the empty word as coaugmentation.
Define the unary and binary operations by
\begin{equation}
\label{eq:bar-corestrictions}
b_1(ta)=-t(d_Ha),\qquad
b_\mu(ta\otimes tb)=(-1)^{|a|}t(ab).
\end{equation}
Extend them as coderivations.
Their degree is one.
For homogeneous inputs,
\begin{equation}
\label{eq:bar-binary-sign}
b_\mu[ a_1|\cdots|a_n]
=\sum_{i=1}^{n-1}(-1)^{\sum_{j\leq i}|a_j|-i+1}
 [a_1|\cdots|a_i a_{i+1}|\cdots|a_n],
\end{equation}
where the brackets denote $ta_1\otimes\cdots\otimes ta_n$.
This follows by moving a degree-one coderivation past the first $i-1$ shifted inputs.

For three inputs, the two binary contractions give
\begin{equation}
\label{eq:associative-two-edge}
b_\mu^2[a|b|c]
=(-1)^{|b|}\,t\bigl((ab)c-a(bc)\bigr)=0.
\end{equation}
This is the first associativity condition in the bar differential.
The algebraic bar construction is also described in
Loday--Vallette, Section~2.2.1~\cite{LV12}.
Equations~\eqref{eq:bar-corestrictions}--\eqref{eq:bar-binary-sign} fix the cohomological signs used here.

\begin{proposition}[Determinant comparison]
\label{prop:bar-determinant-comparison}
Let $L_n=\bigwedge^n k^{[n]}[n]$, with generator $\omega_n=e_1\wedge\cdots\wedge e_n$.
The degree-zero isomorphism
\begin{equation}
\label{eq:bar-determinant-map}
\Phi_n[a_1|\cdots|a_n]
=(-1)^{\sum_{j=1}^n(n-j)|a_j|}
 \omega_n\otimes a_1\otimes\cdots\otimes a_n
\end{equation}
identifies the bar operations with
\[
\Phi b_1\Phi^{-1}=(-1)^n d_{\bar H^{\otimes n}},\qquad
\Phi b_\mu\Phi^{-1}=\sum_{i=1}^{n-1}\iota_{e_i}\otimes\mu_i.
\]
Here $\mu_i$ multiplies positions $i,i+1$ and retains the right vertex label.
After deletion, relabel the surviving vertices in increasing order.
The last vertex is never removed.
Its degree-minus-one line supplements the determinant of the $n-1$ internal cuts.
\end{proposition}

\begin{proof}
Moving each $t$ to the left across preceding algebra factors gives the exponent in~\eqref{eq:bar-determinant-map}.
For multiplication at positions $i,i+1$, the difference between the input and output exponents is
$\sum_{j\leq i}|a_j|$ modulo two.
Adding the exponent in~\eqref{eq:bar-binary-sign} leaves $i-1$ modulo two.
This is precisely the determinant contraction sign.
For the unary operation at position $i$, the exponent in $\Phi_n$ changes by $n-i$.
The unary coderivation contributes
$1+\sum_{j<i}(|a_j|-1)$.
Their sum is $n+\sum_{j<i}|a_j|$ modulo two.
This is the sign in $(-1)^n d_{\bar H^{\otimes n}}$.
\end{proof}

The coefficient maps $\mu_i$ commute on disjoint cuts.
On adjacent cuts their two composites agree by associativity.
Deleting cuts therefore gives a commuting coefficient diagram on every ordered word.
Theorem~\ref{thm:exact-coefficient-criterion}, with the last vertex line retained, proves
$(b_1+b_\mu)^2=0$ in every length.
This also proves compatibility with internal degrees and the augmentation ideal.

Now let $\nu:\bar H\otimes\bar H\to\bar H$ be another degree-zero cochain map.
Use the same sign as~\eqref{eq:bar-corestrictions} to define $b_\nu$.

\begin{theorem}[Two associative bar contractions]
\label{thm:two-associative-contractions}
Suppose $\nu$ is associative.
Then $b_\mu^2=b_\nu^2=0$, and $b_\mu b_\nu+b_\nu b_\mu=0$ if and only if
\begin{align}
\label{eq:compatible-products}
\mu(\nu(a,b),c)+\nu(\mu(a,b),c)
=\mu(a,\nu(b,c))+\nu(a,\mu(b,c))
\end{align}
for every $a,b,c\in\bar H$.
Equivalently, $\mu+z\nu$ is associative over $k[z]$.
Here $z$ is a formal variable of degree zero.
\end{theorem}

\begin{proof}
On three homogeneous inputs, the mixed operator equals $(-1)^{|b|}t$ times the difference of the two sides of~\eqref{eq:compatible-products}.
Thus anticommutation implies that identity.
For longer words, operations on disjoint pairs have equal coefficient composites and opposite determinant signs.
Operations on overlapping pairs occupy three consecutive positions.
Their four contributions are the same identity, tensored with the unchanged inputs.
Thus the identity implies anticommutation in every length.
Expanding associativity of $\mu+z\nu$ gives associativity of $\mu$, equation~\eqref{eq:compatible-products}, and associativity of $\nu$ as its three coefficients.
Coefficient comparison in $k[z]$ requires no hypothesis on the characteristic.
\end{proof}

\begin{example}[Separate associativity is insufficient]
\label{ex:incompatible-products}
Let $\bar H=kx\oplus ky$ in degree zero.
Define $\mu(x,x)=x$ and $\nu(y,y)=x$, with every other basis product zero.
Both products are associative.
For $\mu$, only an all-$x$ triple can have a nonzero iterated product.
Every iterated $\nu$-product vanishes.
For $(a,b,c)=(y,y,x)$, the left side of~\eqref{eq:compatible-products} is $x$ and the right side is zero.
Thus the mixed bar anticommutator is nonzero whenever $k$ is nonzero.
Adjoining a unit to $(\bar H,\mu)$ gives an augmented algebra with this augmentation ideal.
\end{example}

\begin{example}[A noncommutative family]
\label{ex:inserted-element}
Choose a closed element $h\in H^0$ and define $\nu(a,b)=ahb$.
The augmentation ideal contains $ahb$ whenever $a,b\in\bar H$.
The map is a cochain map because $d_Hh=0$.
Associativity gives
\[
(ahb)hc=ah(bhc),\qquad
(ahb)c+(ab)hc=a(bhc)+ah(bc).
\]
Thus the two bar contractions anticommute.
No centrality assumption on $h$ is needed.
For $H=k\langle x,y\rangle$, with zero differential and its usual augmentation, take $h=x$.
The products $yxyy$ and $yyxy$ can occur in separate mixed paths and are distinct.
Their four-term sum still cancels.
This example shows why requiring every mixed coefficient square to commute would be unnecessarily restrictive.
\end{example}

\subsection{Logarithmic residues and the de Rham differential}
\label{subsec:logarithmic-comparison}

Let $Y$ be a smooth complex manifold and let $D=\bigcup_{e\in E}D_e$ be a simple normal crossing divisor.
The components are labeled, and $E$ is finite.
For $S\subseteq E$, write $D_S=\bigcap_{e\in S}D_e$, with $D_\varnothing=Y$.
Let $D^S$ be the divisor on $D_S$ cut out by the remaining components.
Empty strata contribute zero.
All statements are local statements of sheaves on $Y$, using direct image from each stratum.

On a coordinate chart with $D_e=\{x_e=0\}$, define $r_{S,e}$ by~\eqref{eq:left-residue}.
Its domain and target are
\[
r_{S,e}:\Omega^q_{D_S}(\log D^S)
\longrightarrow\Omega^{q-1}_{D_{S+e}}(\log D^{S+e}).
\]
Changing $x_e$ to $u x_e$, with $u$ invertible, adds the regular form $du/u$ to $dx_e/x_e$.
It therefore leaves the residue unchanged.
Restriction of the coefficient to $D_{S+e}$ makes the formula independent of its representative.
This constructs the map on overlapping coordinate charts.

\begin{proposition}[Residue identities]
\label{prop:residue-identities}
With the left-normal convention, the maps satisfy
\begin{equation}
\label{eq:residue-relations}
r_{S+e,f}r_{S,e}=-r_{S+f,e}r_{S,f},\qquad
r_{S,e}d_{D_S}=-d_{D_{S+e}}r_{S,e}.
\end{equation}
The first equation uses distinct $e,f\notin S$.
Consequently, on
\[
\mathcal R=\bigoplus_{S\subseteq E}
\Omega^\bullet_{D_S}(\log D^S)[-2|S|],
\]
the sum $d_{\mathrm{dR}}+R$, with $R=\sum r_{S,e}$, has degree one and squares to zero.
A partition $E=E_a\sqcup E_b$ gives $R=R_a+R_b$ with
$R_a^2=R_b^2=R_aR_b+R_bR_a=0$.
\end{proposition}

\begin{proof}
Only terms containing both $d\log x_e$ and $d\log x_f$ contribute to a double residue.
Moving these two forms to the front in opposite orders changes the sign once.
Coefficient restriction to their common zero locus is independent of order.
This proves the first equation.
For the second, write $\alpha=d\log x_e\wedge\beta+\gamma$.
The Leibniz rule gives
$d\alpha=-d\log x_e\wedge d\beta+d\gamma$.
The term $d\gamma$ has no logarithmic pole along $D_e$.
Restriction of $d\beta$ commutes with exterior differentiation.
Therefore $r_{S,e}d\alpha=-d(r_{S,e}\alpha)$.
This agrees with the shifted residue sequence in
Stacks Project, Lemma~50.15.2~\cite{StacksResidue}.

A $q$-form on $D_S$ has total degree $q+2|S|$ in $\mathcal R$.
A residue lowers $q$ by one and raises $|S|$ by one.
Its total degree is therefore one.
The shifts are even, so the internal differential remains $d_{\mathrm{dR}}$.
The identities just proved cancel every double residue and every mixed term with $d_{\mathrm{dR}}$.
The color decomposition follows by grouping pairs of distinct components.
\end{proof}

To identify the determinant explicitly, take a product chart
$U=\Delta^E\times Z$, with coordinates $x_e$ on the normal discs.
For $S\subseteq E$, write $U_S=\{x_e=0:e\in S\}\subset U$ and
let $p_S:U_S\to Z$ be the projection.
Work with complexes of complex vector spaces on these fixed charts.
Put $M^\bullet=\Gamma(Z,\Omega_Z^\bullet)$, with its holomorphic de Rham differential $d_Z$.
Choose an order on $E$ and use the inherited order on every $F=E\setminus S$.
Let $\lambda_F$ be the wedge of the $d\log x_e$, $e\in F$, in that order.
The empty wedge is $1$.

Define the direct sum over every contracted set:
\[
C_E(Z)=\bigoplus_{S\subseteq E}L(E\setminus S)\otimes M^\bullet.
\]
On its $S$-summand, set
\begin{equation}
\label{eq:normal-log-total-differential}
D_E(\omega_F\otimes\beta)
=(-1)^{|F|}\omega_F\otimes d_Z\beta
+\sum_{e\in F}\iota_e\omega_F\otimes\beta.
\end{equation}
The first term stays in the $S$-summand.
The term indexed by $e$ lies in the $S\cup\{e\}$-summand.
Thus $D_E=\delta+\sum_e\iota_e\otimes\mathrm{id}$, with each identity map acting between the corresponding copies of $M^\bullet$.
The coefficient squares commute, so Theorem~\ref{thm:exact-coefficient-criterion} gives $D_E^2=0$.

Inside the complex $\mathcal R(U)$ of sections of the residue sheaves, put
\[
N_E(Z)=\bigoplus_{S\subseteq E}
\lambda_{E\setminus S}\wedge p_S^*M^\bullet.
\]
Here the $S$-summand is on $U_S$.
It has the grading inherited from $\mathcal R$:
$\lambda_F\wedge p_S^*\beta$ has total degree $|\beta|+|F|+2|S|$.
The direct-sum map
\begin{equation}
\label{eq:log-symbol-comparison}
\Psi:C_E(Z)[-2|E|]\longrightarrow N_E(Z)\subset\mathcal R(U),
\qquad
\omega_F\otimes\beta\longmapsto\lambda_F\wedge p_S^*\beta
\end{equation}
is a chain isomorphism onto this subcomplex.

Indeed, its source degree is $|\beta|-|F|+2|E|$.
Since $|E|=|F|+|S|$, this equals the target degree.
The global shift is even, so it leaves the differential $D_E$ unchanged.
The normal logarithmic forms are closed, and hence
\[
d_{\mathrm{dR}}(\lambda_F\wedge p_S^*\beta)
=(-1)^{|F|}\lambda_F\wedge p_S^*(d_Z\beta).
\]
This term lies in the same summand and agrees with the first term of~\eqref{eq:normal-log-total-differential}.
For $e=e_i$ in the ordered set $F$, the residue is
\[
r_{S,e}(\lambda_F\wedge p_S^*\beta)
=(-1)^{i-1}\lambda_{F\setminus\{e\}}\wedge p_{S\cup\{e\}}^*\beta.
\]
This lies in the included $S\cup\{e\}$-summand and agrees with the corresponding determinant contraction.
These two formulas prove closure of $N_E(Z)$ under $d_{\mathrm{dR}}+R$ and prove
$(d_{\mathrm{dR}}+R)\Psi=\Psi D_E$.
On each summand, the normal logarithmic wedge has a unique coefficient pulled back from $Z$.
Extracting this coefficient gives the inverse to $\Psi$ on its image.

For $E=\{e\}$ and $Z$ a point, the direct sum has two terms.
Under $\Psi$, its differential is $d\log x_e\mapsto1$ on $U_{\{e\}}$.
Their total degrees are one and two, respectively.
The first term has $S=\varnothing$, and the second has $S=\{e\}$.

This comparison concerns the specified direct sum of normal logarithmic forms on the product strata.
General logarithmic forms also have regular normal directions and variable coefficients.
A determinant line alone is not the full sheaf of logarithmic forms.
The global residue identities hold on that full sheaf by Proposition~\ref{prop:residue-identities}.

\subsection{The coefficient obligation for chiral contractions}
\label{subsec:chiral-coefficient-obligation}

In a chiral construction, a contraction can include a collision map on states,
a residue, and a pushforward from a boundary stratum.
Each composite must have a specified target after the first collision.
A second contraction acts on that target, which contains the merged state.

Theorem~\ref{thm:exact-coefficient-criterion} applies after these coefficient complexes and degree-zero maps have been constructed.
It determines the orientation sign and the exact remaining coefficient identity.
For disjoint fixed-color contractions, equality of the two unsigned composites suffices.
When both kinds of contraction can remove either edge, equation~\eqref{eq:coefficient-defects} gives the required four-term identity.
The associative construction proves this identity on its stated tensor-coalgebra carrier.
The logarithmic construction proves it for the stated residue sheaves.
A comparison with a chiral bar complex must intertwine its native collision maps with these maps.
Neither a determinant line nor an equality of scalar residues constructs that comparison.

The de Rham differential is a distinct degree-one operation.
With left-normal residues, its anticommutation follows from the second identity in~\eqref{eq:residue-relations}.
Here $d(d\log x_e)=0$ as a meromorphic differential form.
No current supported on a divisor occurs in this de Rham complex.
Replacing meromorphic forms by currents changes the carrier and requires its own comparison.
